\documentclass[12pt]{article}

\usepackage[margin=1in]{geometry}
\usepackage{amsmath, amssymb, amsthm}
\usepackage{setspace}
\usepackage{hyperref}

\setlength{\parskip}{1em}
\setlength{\parindent}{0pt}
\onehalfspacing

\title{Induction Proofs for the Sum of the First \(n\) Integers and the Sum of the First \(n\) Cubes}
\author{}
\date{}

\begin{document}

\maketitle

\section*{Historical Note}

Carl Friedrich Gauss (1777--1855) is widely regarded as one of the great mathematicians of all time, with major contributions to number theory, geometry, astronomy, geodesy, and physics.\footnote{Encyclopaedia Britannica, ``Carl Friedrich Gauss,'' accessed 2026.} A famous classroom story says that when Gauss was a child, his teacher asked the class to add the integers from \(1\) to \(100\), expecting the exercise to take a while. Gauss allegedly saw immediately that the numbers could be paired as
\[
1+100,\quad 2+99,\quad 3+98,\quad \dots
\]
so that each pair sums to \(101\), and since there are \(50\) such pairs, the total is
\[
50\cdot 101 = 5050.
\]
That gives the general idea behind the formula
\[
1+2+\cdots+n = \frac{n(n+1)}{2}.
\]

This story is repeated in many biographies of Gauss, and the pairing idea is the standard explanation associated with his name.\footnote{MacTutor History of Mathematics, ``Carl Friedrich Gauss,'' University of St Andrews, accessed 2026.} Historians have also noted that the exact schoolroom version of the story appears in multiple later retellings, so it is safest to treat it as a famous anecdote rather than a perfectly documented transcript of a single classroom moment.\footnote{Discussion summarizing the source history of the anecdote: \textit{Did Gauss find the formula for \(1+2+\cdots+n\) in elementary school?}, History of Science and Mathematics Stack Exchange, accessed 2026.}

\vspace{4cm}

\newpage

\section*{Problem Statement}

Let
\[
S(n)=1+2+\cdots+n
\]
be the sum of the first \(n\) natural numbers, and let
\[
C(n)=1^3+2^3+\cdots+n^3
\]
be the sum of the first \(n\) cubes.

Prove by induction on \(n\) that

\[
\text{(a)}\quad S(n)=\frac{1}{2}n(n+1)
\]

and

\[
\text{(b)}\quad C(n)=\frac{1}{4}(n^4+2n^3+n^2)=\frac{1}{4}n^2(n+1)^2.
\]

Then conclude that

\[
C(n)=S(n)^2.
\]

\vspace{6cm}

\newpage

\section*{Part (a): Prove \(S(n)=\frac{1}{2}n(n+1)\) by Induction}

\subsection*{Base Case - a.k.a. P(1) case from discrete}

For \(n=1\),
\[
S(1)=1.
\]
The formula gives
\[
\frac{1}{2}(1)(1+1)=\frac{1}{2}(1)(2)=1.
\]
So the statement is true for \(n=1\).

\vspace{3cm}

\subsection*{Inductive Hypothesis}

Assume that for some \(k\geq 1\),
\[
S(k)=\frac{1}{2}k(k+1).
\]

\vspace{3cm}

\subsection*{Inductive Step}

We must show that
\[
S(k+1)=\frac{1}{2}(k+1)(k+2).
\]

Starting from the definition,
\[
S(k+1)=S(k)+(k+1).
\]
Using the inductive hypothesis,
\[
S(k+1)=\frac{1}{2}k(k+1)+(k+1).
\]
Factor out \((k+1)\):
\[
S(k+1)=(k+1)\left(\frac{1}{2}k+1\right).
\]
Write \(1\) as \(\frac{2}{2}\):
\[
S(k+1)=(k+1)\left(\frac{1}{2}k+\frac{2}{2}\right).
\]
So
\[
S(k+1)=(k+1)\left(\frac{k+2}{2}\right)
=\frac{1}{2}(k+1)(k+2).
\]

Therefore, by induction,
\[
S(n)=\frac{1}{2}n(n+1)
\]
for all natural numbers \(n\).

\vspace{5cm}

\newpage

\section*{Part (b): Prove \(C(n)=\frac{1}{4}n^2(n+1)^2\) by Induction}

\subsection*{Base Case}

For \(n=1\),
\[
C(1)=1^3=1.
\]
The formula gives
\[
\frac{1}{4}(1)^2(1+1)^2=\frac{1}{4}(1)(4)=1.
\]
So the statement is true for \(n=1\).

\vspace{3cm}

\subsection*{Inductive Hypothesis}

Assume that for some \(k\geq 1\),
\[
C(k)=\frac{1}{4}k^2(k+1)^2.
\]

\vspace{3cm}

\subsection*{Inductive Step}

We must show that
\[
C(k+1)=\frac{1}{4}(k+1)^2(k+2)^2.
\]

Starting from the definition,
\[
C(k+1)=C(k)+(k+1)^3.
\]
Using the inductive hypothesis,
\[
C(k+1)=\frac{1}{4}k^2(k+1)^2+(k+1)^3.
\]
Factor out \((k+1)^2\):
\[
C(k+1)=(k+1)^2\left(\frac{1}{4}k^2+(k+1)\right).
\]
Put everything over a common denominator:
\[
C(k+1)=(k+1)^2\left(\frac{k^2+4k+4}{4}\right).
\]
Since
\[
k^2+4k+4=(k+2)^2,
\]
we get
\[
C(k+1)=(k+1)^2\left(\frac{(k+2)^2}{4}\right).
\]
Therefore,
\[
C(k+1)=\frac{1}{4}(k+1)^2(k+2)^2.
\]

So by induction,
\[
C(n)=\frac{1}{4}n^2(n+1)^2
\]
for all natural numbers \(n\).

\vspace{5cm}

\newpage

\section*{Conclusion}

From part (a),
\[
S(n)=\frac{1}{2}n(n+1).
\]
Squaring both sides gives
\[
S(n)^2=\left(\frac{1}{2}n(n+1)\right)^2
=\frac{1}{4}n^2(n+1)^2.
\]

From part (b),
\[
C(n)=\frac{1}{4}n^2(n+1)^2.
\]

Therefore,
\[
C(n)=S(n)^2.
\]

So the sum of the first \(n\) cubes is equal to the square of the sum of the first \(n\) natural numbers:
\[
1^3+2^3+\cdots+n^3 = \left(1+2+\cdots+n\right)^2.
\]

\vspace{6cm}

\newpage

\section*{Short Slide-Friendly Version}

\textbf{Key idea behind the sum \(1+2+\cdots+n\):}

Pair the first and last terms:
\[
1+n,\quad 2+(n-1),\quad 3+(n-2),\dots
\]
Each pair sums to \(n+1\), which leads to
\[
S(n)=\frac{n(n+1)}{2}.
\]

\vspace{3cm}

\textbf{Key induction result for cubes:}
\[
1^3+2^3+\cdots+n^3=\frac{1}{4}n^2(n+1)^2.
\]

\vspace{3cm}

\textbf{Beautiful final identity:}
\[
1^3+2^3+\cdots+n^3
=
\left(1+2+\cdots+n\right)^2.
\]

\vspace{3cm}

\textbf{Interpretation:}

The sum of cubes is not just another polynomial formula. It is exactly the square of the triangular number.

\vfill

\section*{Why This Should Look Familiar from Calculus}

These formulas also show up later in calculus when definite integrals are built from Riemann sums. In that setting, an area is approximated by adding many small rectangle areas together, and the algebra often produces sums such as
\[
1+2+\cdots+n,\qquad 1^2+2^2+\cdots+n^2,\qquad 1^3+2^3+\cdots+n^3.
\]
For example, when finding an integral such as
\[
\int_0^1 x^2\,dx
\]
from first principles, the Riemann sum expands into a constant multiple of
\[
1^2+2^2+\cdots+n^2.
\]
Likewise, higher-degree polynomial integrals produce higher-power sums. That is one reason these summation formulas matter: they are not just number patterns, but part of the machinery used to pass from finite sums to exact area under a curve.

\section*{References}

\begin{enumerate}
    \item Encyclopaedia Britannica, ``Carl Friedrich Gauss,'' accessed 2026.
    \item MacTutor History of Mathematics Archive, University of St Andrews, ``Carl Friedrich Gauss,'' accessed 2026.
    \item History of Science and Mathematics Stack Exchange, discussion on the source history of the Gauss schoolroom anecdote, accessed 2026.
\end{enumerate}

\end{document}
